% This is samplepaper.tex, a sample chapter demonstrating the
% LLNCS macro package for Springer Computer Science proceedings;
% Version 2.21 of 2022/01/12
%
\documentclass[runningheads]{llncs}
%
\usepackage[T1]{fontenc}
% T1 fonts will be used to generate the final print and online PDFs,
% so please use T1 fonts in your manuscript whenever possible.
% Other font encondings may result in incorrect characters.
%
\usepackage{graphicx}
% Used for displaying a sample figure. If possible, figure files should
% be included in EPS format.
%
\usepackage[breaklinks]{hyperref}
\usepackage{xurl}

\usepackage[acronym]{glossaries}
\usepackage{listings}
%% auto break lines
\lstset{
  breaklines=true,
  basicstyle=\fontsize{8}{9}\selectfont\ttfamily,
}
\usepackage[style=numeric]{biblatex}
\addbibresource{references.bib}

\usepackage{subfiles} % Best loaded last in the preamble


\newacronym[longplural=Knowledge Graphs]{kg}{KG}{Knowledge Graph}
\newacronym{dkg}{DKG}{Decentralized Knowledge Graph}
\newacronym{sssom}{SSSOM}{Simple Standard for Sharing Ontological Mappings}
\newacronym{iri}{IRI}{Internationalized Resource Identifier}
\newacronym{ltqp}{LTQP}{Link Traversal Query Processing}
\newacronym{rdf}{RDF}{Resource Description Framework}
\newacronym{http}{HTTP}{Hypertext Transfer Protocol}
\newacronym{https}{HTTPS}{Hypertext Transfer Protocol Secure}

%%
%% end of the preamble, start of the body of the document source.
\begin{document}

%%
%% The "title" command has an optional parameter,
%% allowing the author to define a "short title" to be used in page headers.
\title{Demonstrating Online Schema Alignment in Decentralized Knowledge Graph Querying}
\titlerunning{Demonstrating Online Schema Alignment}
%
%\titlerunning{Abbreviated paper title}
% If the paper title is too long for the running head, you can set
% an abbreviated paper title here
%
\author{Bryan-Elliott Tam\inst{1}\orcidID{0000-0003-3467-9755} \and
Pieter Colpaert\inst{1}\orcidID{0000-0001-6917-2167} \and
Ruben Taelman\inst{1}\orcidID{0000-0001-5118-256X}}
%
\authorrunning{B-E. Tam et al.}
% First names are abbreviated in the running head.
% If there are more than two authors, 'et al.' is used.
%
\institute{
  IDLab,
  Dept.\ Electronics and Information Systems,
  Ghent University -- imec, Belgium
  %\email{\{firstname.lastname\}@ugent.be}
}

\maketitle

\subfile{section/abstract}

%%
%% Keywords. The author(s) should pick words that accurately describe
%% the work being presented. Separate the keywords with commas.
\keywords{SPARQL \and Decentralization \and Link Traversal \and Schema Alignment}

\subfile{section/introduction}
% \subfile{section/preliminary}
\subfile{section/method}
\subfile{section/demo}
\subfile{section/conclusion}
%%
%% The next two lines define the bibliography style to be used, and
%% the bibliography file.
%\bibliographystyle{ACM-Reference-Format}
%\bibliography{}

\printbibliography

\end{document}

%%
%% End of file `sample-sigconf.tex'.
